\documentclass[12pt]{article}
\usepackage{amsmath,amssymb}

\begin{document}
\section*{{2024 AMC 12B Problems/Problem 23}}
Source: AoPS Wiki

\subsection*{Solution}
To find the height of the pyramid, we need the length from the center of the octagon (denote as $I$ ) to its vertices and the length of AV. From symmetry, we know that $\overline{AV} = \overline{DV}$ , therefore $\triangle{AVD}$ is a 45-45-90 triangle. Denote $\overline{AV}$ as $x$ so that $\overline{AD} = x\sqrt{2}$ . Doing some geometry on the isosceles trapezoid $ABCD$ (we know this from the fact that it is a regular octagon) reveals that $\overline{AD}=1+2(\sqrt{2}/2)=1+\sqrt{2}$ and $\overline{AV}=(\overline{AD})/\sqrt{2}=(\sqrt{2}+2)/2$ . To find the length $\overline{IA}$ , we cut the octagon into 8 triangles, each with a smallest angle of 45 degrees. Using the law of cosines on $\triangle{AIB}$ we find that ${\overline{IA}}^2=(2+\sqrt{2})/2$ . Finally, using the pythagorean theorem, we can find that ${\overline{IV}}^2={\overline{AV}}^2-{\overline{IA}}^2= {((\sqrt{2}+2)/2)}^2 - (2+\sqrt{2})/2 = \boxed{(1+\sqrt{2})/2}$ which is answer choice $\boxed{B}$ . ～username2333
~hashbrown2009

\end{document}
